%%% FIELD proof-prop-one-path-limit
Put \(h=1/8\).  For the policy grid in \(\textnormal{def:discontinuity-witness}\), the six latent-index bins, in increasing order, have lengths
\[
(h,h,1/4,1/4,h,h).
\]
The first two bins are always-treated, the middle two are complier bins, and the last two are never-treated.  Fix \(t\in(0,1/4)\).  We first derive the unique feasible path and its complete policy image from the allocation representation used in \(\textnormal{thm:continuity-frontier}\).

Index a path node by the realized pair \((Y_0,Y_1)\).  Let \(m_{ij}\) be the allocation at \((y_i,y_j)\), aggregated over the two complier bins.  The zero second coordinates of both complier margins imply, by nonnegativity, that every \(m_{2j}\) and every \(m_{i2}\) is zero.  The remaining row and column balance equations include
\[
m_{11}+m_{13}=\frac14,
\qquad
m_{11}+m_{31}=\frac14-t.
\]
Consequently
\[
m_{13}-m_{31}=t>0,
\]
so \(m_{13}>0\).  A coordinatewise monotone path containing \((1,3)\) cannot also contain a node \((i,j)\) with \(i>1\) and \(j<3\); hence \(m_{31}=0\).  The balance equations then force
\[
m_{11}=\frac14-t,
\qquad
m_{13}=t,
\qquad
m_{33}=\frac14.
\]
Because all three displayed masses are strictly positive, a feasible maximal path must pass successively through \((1,1)\), \((1,3)\), and \((3,3)\).  There is exactly one such path: it makes both increments in the second coordinate before both increments in the first coordinate.  Thus the other five paths have empty allocation fibers and this path is the unique survivor.

On the surviving path, the always-treated margin \(a_1=(0,1/4,0)\) forces mass \(h\) at \((0,1)\) in each of the first two bins.  Similarly, the never-treated margin \(n_0=(0,1/4,0)\) forces mass \(h\) at \((1,2)\) in each of the last two bins.  Let \(v\) be the mass assigned to \((0,2)\) in the first complier bin.  Nonnegativity and the total mass \(t\) at that node give \(0\le v\le t\).  Conversely, every \(v\in[0,t]\) is feasible: in the two complier bins use respectively
\[
\begin{array}{c|ccc}
 & (0,0)&(0,2)&(2,2)\\ \hline
 b_3&1/4-t&v&t-v\\
 b_4&0&t-v&1/4-t+v.
\end{array}
\]
All entries are nonnegative, each row sums to \(1/4\), and the three column sums are \(1/4-t\), \(t\), and \(1/4\), respectively.  An arbitrary feasible allocation may split the first and third column masses differently between the two complier bins, but its policy means depend on that split only through \(v\).  Indeed, writing the three policy means as \(\mu_1,\mu_2,\mu_3\), the balance equations give
\[
\begin{aligned}
\mu_1
 &=h+2\left(\frac14\right)+\frac14
 =\frac78,\\
\mu_2
 &=\frac14+2\left(v+\frac14\right)+\frac14
 =1+2v,\\
\mu_3
 &=\frac14+2\left(\frac14+t\right)+2h+h
 =\frac98+2t.
\end{aligned}
\]
Here the first equality uses treatment only in the first bin; the second uses treatment in the first complier bin and the identity that total mass at \((2,2)\) is \(1/4\); and the third uses treatment through the first never-treated bin.  The allocation representation is sharp, so the feasible construction and the preceding necessity argument establish
\[
S_t=\left\{\left(\frac78,r,\frac98+2t\right):1\le r\le1+2t\right\}
\qquad (0<t<1/4).
\tag{1}
\]
In particular, the same component survives throughout a neighborhood of every fixed \(t_0\in(0,1/4)\); the survival condition of \(\textnormal{thm:continuity-frontier}\) therefore also verifies Hausdorff continuity there.

We next compute the local distance exactly.  Suppose \(s\ge t>0\).  The middle-coordinate interval for \(S_t\) is contained in that for \(S_s\), while their constant third coordinates differ by \(2(s-t)\).  Hence every point of \(S_t\) has distance exactly \(2(s-t)\) from the point of \(S_s\) with the same middle coordinate.  In the reverse direction, match \(r\in[1,1+2s]\) to \(\min\{r,1+2t\}\).  Both the middle-coordinate discrepancy and the third-coordinate discrepancy are at most \(2(s-t)\), and the latter is always equal to \(2(s-t)\).  Thus both directed excesses equal \(2(s-t)\).  Interchanging \(s\) and \(t\) yields
\[
d_H(S_s,S_t)=2|s-t|
\tag{2}
\]
whenever \(s,t>0\), and hence locally around \(t_0\).

It remains to derive the sampling limit.  Let \(N_1=\sum_{i=1}^n\mathbf 1\{Z_i=1\}\).  For an observation in arm \(Z=1\), define \(X_i=1/2\) on the treated cell with outcome \(2\), \(X_i=-1/2\) on the treated cell with outcome \(0\), and \(X_i=0\) on the other cells.  On \(\{N_1>0\}\), the unclipped contrast is
\[
\widetilde t
=\frac{\widehat q_1(1,3)-\widehat q_1(1,1)}2
=\frac1{N_1}\sum_{i=1}^n\mathbf 1\{Z_i=1\}X_i.
\]
Under \(q_{t_0}\), the two changing conditional arm probabilities are \(q_{t_0,1}(1,3)=1/4+t_0\) and \(q_{t_0,1}(1,1)=1/4-t_0\).  Therefore
\[
\mathbb E[X_i\mid Z_i=1]=t_0,
\qquad
\mathbb E[X_i^2\mid Z_i=1]
=\frac14\left\{\left(\frac14+t_0\right)+\left(\frac14-t_0\right)\right\}
=\frac18,
\]
and
\[
\operatorname{Var}(X_i\mid Z_i=1)=\frac18-t_0^2>0.
\tag{3}
\]
Set \(A_i=\mathbf 1\{Z_i=1\}(X_i-t_0)\).  The variables \(A_i\) are bounded and independent, with mean zero and variance
\[
\mathbb E[A_i^2]=\pi_1\left(\frac18-t_0^2\right).
\]
For every fixed real \(u\), boundedness and the Taylor expansion of the exponential give
\[
\mathbb E\exp\left(\frac{iuA_i}{\sqrt n}\right)
=1-\frac{u^2}{2n}\pi_1\left(\frac18-t_0^2\right)+o(n^{-1}).
\]
Raising this expression to the \(n\)-th power proves directly that
\[
\frac1{\sqrt n}\sum_{i=1}^n A_i
\ \Rightarrow\ 
N\!\left(0,\pi_1\left(\frac18-t_0^2\right)\right).
\tag{4}
\]
Also, independence and a variance calculation give \(N_1/n\to\pi_1\) in probability.  Since \(\pi_1>0\), \(\Pr(N_1=0)=(1-\pi_1)^n\to0\), and on \(\{N_1>0\}\),
\[
\sqrt n(\widetilde t-t_0)
=\frac{n}{N_1}\frac1{\sqrt n}\sum_{i=1}^n A_i.
\]
The centered sum in (4) is tight, so replacing \(n/N_1\) by \(1/\pi_1\) changes the last display by a term converging to zero in probability.  Equations (3)--(4), together with \(\pi_1=1/2\), therefore imply
\[
\sqrt n(\widetilde t-t_0)
\ \Rightarrow\ 
N\!\left(0,\frac{1/8-t_0^2}{\pi_1}\right)
=N\!\left(0,\frac14-2t_0^2\right).
\tag{5}
\]
The same calculation gives \(\widetilde t\to t_0\) in probability.  Because \(t_0\) is strictly between \(0\) and \(1/4\), the clipping map equals the identity on a fixed neighborhood of \(t_0\).  Hence \(\Pr(\widehat t=\widetilde t)\to1\), and (5) holds with \(\widehat t\) in place of \(\widetilde t\).  With probability tending to one, \(\widehat t\) and \(t_0\) are in the neighborhood where (2) applies.  Taking absolute values in (5) and using (2) gives
\[
\sqrt n\,d_H(S_{\widehat t},S_{t_0})
\ \Rightarrow\ 
2\sqrt{\frac14-2t_0^2}\,|N(0,1)|.
\]

Finally, let \(c_{1-\alpha}\) denote the \((1-\alpha)\)-quantile of \(|N(0,1)|\), and put \(\widehat\sigma=\sqrt{1/4-2\widehat t^{,2}}\).  Continuity and consistency imply \(\widehat\sigma\to\sqrt{1/4-2t_0^2}>0\) in probability.  Consequently the Hausdorff ball centered at \(S_{\widehat t}\) with radius \(2\widehat\sigma c_{1-\alpha}/\sqrt n\) satisfies
\[
\Pr_{q_{t_0}}\!\left\{d_H(S_{\widehat t},S_{t_0})
\le \frac{2\widehat\sigma c_{1-\alpha}}{\sqrt n}\right\}
\longrightarrow 1-\alpha.
\]
This is the asserted pointwise first-order plug-in folded-normal coverage.
